\documentclass[../../main.tex]{subfiles}

\begin{document}

\section{拓扑群: 一般线性群}\label{section:ch22-拓扑群}

考虑群 $\mathcal{G}$ 以及 $\mathcal{G}$ 上的 Hausdorff 拓扑. 对 $\mathcal{G}$ 的两个成员 $g_{1}$ 与 $g_{2}$, 记群运算为 $g_{1}\cdot g_{2}$, 记群的成员 $g$ 的逆为 $g^{-1}$, 且令 $e$ 为群的单位元. 我们称 $\mathcal{G}$ 是拓扑群, 若从 $\mathcal{G}\times\mathcal{G}$ 到 $\mathcal{G}$ 的映射 $(g_{1},g_{2})\mapsto g_{1}\cdot g_{2}$ 是连续的, 其中 $\mathcal{G}\times\mathcal{G}$ 具有乘积拓扑, 且从 $\mathcal{G}$ 到 $\mathcal{G}$ 的映射 $g\mapsto g^{-1}$ 是连续的. 谈到紧群我们意味着拓扑群作为拓扑空间是紧的. 对 $\mathcal{G}$ 的子集 $\mathcal{G}_{1}$ 和 $\mathcal{G}_{2}$, 我们定义 $\mathcal{G}_{1}\cdot\mathcal{G}_{2}=\{g_{1}\cdot g_{2}\mid g_{1}\in\mathcal{G}_{1},\ g_{2}\in\mathcal{G}_{2}\}$ 与 $\mathcal{G}_{1}^{-1}=\{g^{-1}\mid g\in\mathcal{G}_{1}\}$. 若 $\mathcal{G}_{1}$ 只有一个成员 $g$, 我们记 $\{g\}\cdot\mathcal{G}_{2}$ 为 $g\cdot\mathcal{G}_{2}$.

令 $E$ 为 Banach 空间, 而 $\mathcal{L}(E)$ 为 $E$ 上的连续线性算子的 Banach 空间. $\mathcal{L}(E)$ 中的两个算子的复合也属于 $\mathcal{L}(E)$, 且显然对算子 $T,\ S\in\mathcal{L}(E)$,
\begin{align}
\|S\circ T\|\leqslant\|S\|\cdot\|T\|.
\label{align:ch22-compound-norm}
\end{align}
定义 $\mathrm{GL}(E)$ 为 $\mathcal{L}(E)$ 中的可逆算子族. $\mathcal{L}(E)$ 中的算子是可逆的当且仅当它是一对一与映上的; 根据开映射定理, 它的逆是连续的. 观察到对 $T,\ S\in\mathrm{GL}(E)$, $(S\circ T)^{-1}=T^{-1}\circ S^{-1}$. 因此, 在复合的运算下, $\mathrm{GL}(E)$ 是一个群, 称为 $E$ 的一般线性群. 我们将它的单位元记为 $\mathrm{Id}$. 它也是一个拓扑空间, 其中拓扑由算子范数诱导.

\begin{lemma}[Neumann 级数]\label{lemma:ch22-neumann-series}
令 $E$ 为 Banach 空间而算子 $C\in\mathcal{L}(E)$ 满足 $\|C\|<1$. 则 $\mathrm{Id}-C$ 是可逆的且
\begin{align}
\|(\mathrm{Id}-C)^{-1}\|\leqslant(1-\|C\|)^{-1}.
\label{align:ch22-neumann-estimate}
\end{align}
\end{lemma}

\begin{proof}
我们从\eqref{align:ch22-compound-norm}推出对每个自然数 $k$, $\|C^{k}\|\leqslant\|C\|^{k}$. 因此, 由于 $\|C\|<1$, 实数级数 $\sum_{k=0}^{\infty}\|C^{k}\|$ 收敛. 赋范线性空间 $\mathcal{L}(E)$ 是完备的. 因此算子级数 $\sum_{k=0}^{\infty}C^{k}$ 在 $\mathcal{L}(E)$ 中收敛于连续线性算子. 观察到对所有 $n$,
\begin{align*}
(\mathrm{Id}-C)\circ\left(\sum_{k=0}^{n}C^{k}\right)=\left(\sum_{k=0}^{n}C^{k}\right)\circ(\mathrm{Id}-C)=\mathrm{Id}-C^{n+1}.
\end{align*}
因此级数 $\sum_{k=0}^{\infty}C^{k}$ 收敛到 $\mathrm{Id}-C$ 的逆. 估计式\eqref{align:ch22-neumann-estimate}从 $\mathrm{Id}-C$ 的逆的级数表示得出.
\end{proof}

\begin{theorem}[一般线性群是拓扑群]\label{theorem:ch22-gl-topological}
令 $E$ 为 Banach 空间. 则 $E$ 的一般线性群 $\mathrm{GL}(E)$ 是拓扑群, 群运算是复合, 而拓扑由 $\mathcal{L}(E)$ 上的算子范数诱导.
\end{theorem}

\begin{proof}
对 $\mathrm{GL}(E)$ 中的算子 $T$, $T'$, $S$, $S'$, 观察到
\begin{align*}
T\circ S-T'\circ S'=T\circ(S-S')+(T-T')\circ S'.
\end{align*}
因此, 根据算子范数的三角不等式和不等式\eqref{align:ch22-compound-norm},
\begin{align*}
\|T\circ S-T'\circ S'\|\leqslant\|T\|\cdot\|S-S'\|+\|T-T'\|\cdot\|S'\|.
\end{align*}
复合的连续性从该不等式得到.

若 $S$ 属于 $\mathrm{GL}(E)$ 且 $\|S-\mathrm{Id}\|<1$, 则从等式
\begin{align*}
S^{-1}-\mathrm{Id}=(\mathrm{Id}-S)S^{-1}=(\mathrm{Id}-S)[\mathrm{Id}-(\mathrm{Id}-S)]^{-1}
\end{align*}
以及不等式\eqref{align:ch22-compound-norm}和\eqref{align:ch22-neumann-estimate}, 我们推出
\begin{align}
\|S^{-1}-\mathrm{Id}\|\leqslant\frac{\|S-\mathrm{Id}\|}{1-\|S-\mathrm{Id}\|}.
\label{align:ch22-inverse-near-id}
\end{align}
因此逆在单位元是连续的. 现在令 $T$ 和 $S$ 属于 $\mathrm{GL}(E)$ 且 $\|S-T\|<\|T^{-1}\|^{-1}$. 则
\begin{align*}
\|T^{-1}S-\mathrm{Id}\|=\|T^{-1}(S-T)\|\leqslant\|T^{-1}\|\cdot\|S-T\|<1.
\end{align*}
因此, 若在\eqref{align:ch22-inverse-near-id}中用 $T^{-1}S$ 代替 $S$, 有
\begin{align*}
\|S^{-1}T-\mathrm{Id}\|\leqslant\frac{\|T^{-1}S-\mathrm{Id}\|}{1-\|T^{-1}S-\mathrm{Id}\|}.
\end{align*}
从这个不等式以及等式
\begin{align*}
S^{-1}-T^{-1}=(S^{-1}T-\mathrm{Id})T^{-1}\quad\text{与}\quad T^{-1}S-\mathrm{Id}=T^{-1}(S-T)
\end{align*}
我们推出,
\begin{align*}
\|S^{-1}-T^{-1}\|\leqslant\frac{\|T^{-1}\|^{2}\cdot\|T-S\|}{1-\|T^{-1}\|\cdot\|T-S\|}.
\end{align*}
逆在 $T$ 的连续性从该不等式推出.
\end{proof}

在 $E$ 是欧氏空间 $\mathbf{R}^{n}$ 的情形, 记 $\mathrm{GL}(E)$ 为 $\mathrm{GL}(n,\mathbf{R})$. 若选定了 $\mathbf{R}^{n}$ 的基, 则 $\mathrm{GL}(n,\mathbf{R})$ 上的拓扑满足以下要求: 表示该基的算子的 $n\times n$ 矩阵的每个元素是连续函数.

拓扑群的具有子空间拓扑的子群也是一个拓扑群. 例如, 若 $H$ 是 Hilbert 空间, 则 $\mathrm{GL}(H)$ 的那些在内积下保持不变的算子组成的子集是一个拓扑群, 它称为 $H$ 的正交线性群且记为 $\mathcal{O}(H)$.

\begin{exercise}\label{exercise:ch22-01}
若 $\mathcal{T}$ 是拓扑在 $e$ 的基, 证明 $\{g\cdot\mathcal{O}\mid\mathcal{O}\in\mathcal{T}\}$ 是拓扑在 $g\in\mathcal{G}$ 的基.
\end{exercise}

\begin{exercise}\label{exercise:ch22-02}
证明: 若 $K_{1}$ 和 $K_{2}$ 是 $\mathcal{G}$ 的紧子集, 则 $K_{1}\cdot K_{2}$ 是紧的.
\end{exercise}

\begin{exercise}\label{exercise:ch22-03}
令 $\mathcal{O}$ 为 $e$ 的邻域. 证明存在 $e$ 的邻域 $\mathcal{U}$ 使得 $\mathcal{U}=\mathcal{U}^{-1}$ 且 $\mathcal{U}\cdot\mathcal{U}\subseteq\mathcal{O}$.
\end{exercise}

\begin{exercise}\label{exercise:ch22-04}
证明子群 $H$ 的闭包是 $\mathcal{G}$ 的子群.
\end{exercise}

\begin{exercise}\label{exercise:ch22-05}
令 $\mathcal{G}_{1}$ 和 $\mathcal{G}_{2}$ 为拓扑群而 $h:\mathcal{G}_{1}\to\mathcal{G}_{2}$ 是群同态. 证明 $h$ 是连续的当且仅当它在 $\mathcal{G}_{1}$ 的单位元是连续的.
\end{exercise}

\begin{exercise}\label{exercise:ch22-06}
用压缩映射原理证明引理\cref{lemma:ch22-neumann-series}.
\end{exercise}

\begin{exercise}\label{exercise:ch22-07}
用 $\mathcal{L}(E)$ 的完备性证明: 若 $C\in\mathcal{L}(E)$ 且 $\|C\|<1$, 则 $\sum_{k=0}^{\infty}C^{k}$ 在 $\mathcal{L}(E)$ 中收敛.
\end{exercise}

\begin{exercise}\label{exercise:ch22-08}
证明行列式为 $1$ 的 $n\times n$ 可逆实矩阵的集合是一个拓扑群, 若群运算是矩阵的乘法且拓扑是使得矩阵的每个元素是连续函数的拓扑. 该拓扑群称为特殊线性群, 记为 $\mathrm{SL}(n,\mathbf{R})$.
\end{exercise}

\begin{exercise}\label{exercise:ch22-09}
令 $H$ 为 Hilbert 空间. 证明 $\mathrm{GL}(H)$ 中的算子保持范数当且仅当它保持内积.
\end{exercise}

\begin{exercise}\label{exercise:ch22-10}
考虑具有欧氏内积与范数的 $\mathbf{R}^{n}$. 刻画在规范正交基下表示正交算子的那些 $n\times n$ 矩阵.
\end{exercise}

\begin{exercise}\label{exercise:ch22-11}
证明 $\mathrm{GL}(E)$ 在 $\mathcal{L}(E)$ 中是开的.
\end{exercise}

\begin{exercise}\label{exercise:ch22-12}
证明 $\mathrm{GL}(E)$ 中由恒等算子的线性紧扰动组成的算子集是 $\mathrm{GL}(E)$ 的子群. 它记为 $\mathrm{GL}_{C}(E)$.
\end{exercise}

\end{document}